% titel: Blatt 2, A6: Statik -- Walze im Winkel
% skript: Kap. Statik, Kraeftegleichgewicht

% Krafte auf eine ruhende Walze in einer V-foermigen Rinne.

Eine ruhende Walze in einer Rinne aus zwei schraegen Brettern: die
Summe aller Kraefte (Gewicht, zwei Normalkraefte) ist null. Zerlegt man
geeignet, erhaelt man die Normalkraefte jedes Brettes.

\begin{werkzeugkasten}[Werkzeuge]
\textbf{Statisches Gleichgewicht} (ruhender Koerper):
\begin{equation}
  \sum\vec F=\vec 0\qquad(\vec a=0).
\end{equation}

\textbf{Komponentenzerlegung:} Eine Kraeftegleichung in
$x$- und $y$-Richtung aufspalten. Bei einer schraegen Bretterebene zeigt
die Normalkraft senkrecht zur Brettflaeche, also unter einem
Erg{\"a}nzungs- und Wechselwinkel zur Horizontalen.

Für Bretter unter $\alpha_1$ (links) und $\alpha_2$ (rechts) ergibt sich
aus
\begin{equation}
  N_1\sin\alpha_1=N_2\sin\alpha_2,\qquad
  N_1\cos\alpha_1+N_2\cos\alpha_2=mg
\end{equation}
nach Aufloesen
\begin{equation}
  N_1=\frac{mg\,\sin\alpha_2}{\sin(\alpha_1+\alpha_2)},\qquad
  N_2=\frac{mg\,\sin\alpha_1}{\sin(\alpha_1+\alpha_2)}\;.
\end{equation}
\end{werkzeugkasten}

\begin{werkzeugkasten}[Vorgehen]
\begin{itemize}
  \item Skizze: Gewicht senkrecht nach unten, Normalkr{\"a}fte senkrecht zu
        jeder Brettflaeche.
  \item Kraeftegleichung in zwei unabhaengige Komponenten zerlegen und
        nach $N_1,N_2$ aufloesen.
  \item Die Gegenkraefte der Walze auf die Bretter betragen jeweils
        $-N_1$ bzw.\ $-N_2$ (Actio=Reactio).
  \item Bei ``vierfache Gewichtskraft'' $N_1=N_2=4mg$ setzen und die
        Winkel berechnen.
\end{itemize}
\end{werkzeugkasten}

\begin{achtung}[Typische Fallen]
\begin{itemize}
  \item Die Normalkraft steht \emph{senkrecht} zur Brettflaeche, nicht
        senkrecht zur Horizontalen.
  \item $\sin(\alpha_1+\alpha_2)$ im Nenner; beide Winkel gehen ein.
  \item ``Gegenkraft der Walze auf das Brett'' ist nicht dasselbe wie die
        auf die Walze wirkende Kraft -- Actio/Reactio betrifft
        verschiedene Koerper.
\end{itemize}
\end{achtung}
